\documentclass[12pt,a4paper]{article}
\usepackage[utf8]{inputenc}
\usepackage[turkish]{babel}
\usepackage[T1]{fontenc}
\usepackage{amsmath, amssymb}
\usepackage[hidelinks]{hyperref}
\usepackage{geometry}

\geometry{margin=1in}

% --- Kapak Bilgileri ---
\title{\textbf{Makine Öğrenmesine Giriş}}
\author{Oğuz Taşdemir}
\date{\today}

\begin{document}

\maketitle
\newpage

\tableofcontents
\newpage

% --- Giriş Bölümü ---
\newpage
\section*{Giriş}
\addcontentsline{toc}{section}{Giriş}
Bu çalışma, İstanbul Medeniyet Üniversitesi bünyesinde verilen Makine Öğrenmesine Giriş dersi kapsamında hazırlanan kapsamlı bir ders rehberidir. Temel amaç, dönem boyunca işlenen teorik konuları, algoritmaların çalışma mantığını ve uygulama stratejilerini bir araya getirerek öğrenciye bütünsel bir bakış açısı sunmaktır. 

Döküman; veri hazırlığından derin öğrenmeye kadar uzanan geniş bir yelpazeyi, hem akademik disiplinle hem de anlaşılır bir dille ele almaktadır. Rehberin, sınav hazırlık sürecinde ve projelerde temel bir başvuru kaynağı olması hedeflenmiştir.

\newpage
\section{Makine Öğrenmesinin Temelleri}
Makine öğrenmesi algoritmaları, öğrenme biçimlerine göre üç ana kategoriye ayrılır:

\begin{itemize}
    \item \textbf{Denetimli (Supervised) Öğrenme:} Veride hem girdi hem de hedef etiket (label) bulunur. Model, girdiler ile çıktılar arasındaki ilişkiyi öğrenir (Örn: Sınıflandırma, Regresyon).
    \item \textbf{Denetimsiz (Unsupervised) Öğrenme:} Veride etiket yoktur. Algoritma, verinin kendi içindeki gizli yapıları veya benzerlikleri keşfeder (Örn: Kümeleme, Boyut Düşürme).
    \item \textbf{Pekiştirmeli (Reinforcement) Öğrenme:} Bir ajanın (agent), bir ortamda (environment) en yüksek ödülü almak için deneme-yanılma yoluyla öğrenmesidir.
\end{itemize}

\subsection{Makine Öğrenmesi Nedir?}
Bilgisayarların açıkça programlanmadan veriden öğrenmesini sağlayan bir yapay zeka alt dalıdır. Klasik programlamada girdi ve kurallar verilip çıktı alınırken; makine öğrenmesinde girdi ve çıktılar verilerek bu ikisi arasındaki ilişkiyi temsil eden model (kurallar seti) öğrenilir.

Matematiksel olarak bir öğrenme süreci, bir performans ölçütünü (hata payı gibi) optimize edecek şekilde parametrelerin ($\theta$) güncellenmesidir:
\begin{equation}
    y = g(x | \theta)
\end{equation}
Burada $g$ öğrenilen fonksiyonu, $x$ girdiyi, $\theta$ ise modelin iç parametrelerini temsil eder. Amaç, gerçek değer ile tahmin arasındaki farkı (kayıp/loss) minimize etmektir.

\subsubsection{Denetimli, Denetimsiz ve Yarı-Denetimli Öğrenme}
\begin{itemize}
    \item \textbf{Denetimli Öğrenme (Supervised Learning):} Veri setinde her girdinin bir karşılık gelen etiketi (label) bulunur. İki ana alt başlığa ayrılır:
    \begin{itemize}
        \item \textbf{Sınıflandırma (Classification):} Çıktı kategoriktir (Örn: Spam/Değil, Kredi Riski: Yüksek/Düşük).
        \item \textbf{Regresyon (Regression):} Çıktı süreklidir (Örn: Ev fiyatı tahmini, hisse senedi değeri).
    \end{itemize}
    \item \textbf{Denetimsiz Öğrenme (Unsupervised Learning):} Etiketsiz verilerdeki gizli yapıları ve örüntüleri bulmayı hedefler. Temel uygulamaları \textbf{Kümeleme (Clustering)} ve \textbf{Boyut Düşürme}'dir.
    \item \textbf{Yarı-Denetimli Öğrenme:} Az miktarda etiketli veri ile çok büyük miktarda etiketsiz verinin birlikte kullanıldığı hibrit yöntemdir.
\end{itemize}

\subsubsection{Takviyeli Öğrenme (Reinforcement Learning) Temelleri}
Bir ajanın (agent), bir ortamda (environment) eylemler gerçekleştirerek en yüksek ödülü (reward) toplamaya çalıştığı öğrenme türüdür. "Gecikmiş Ödül" kavramı kritiktir; ajan anlık kazanç yerine uzun vadeli stratejiye odaklanır (Örn: Satranç motorları, otonom sürüş).

\subsubsection{Temel Teoremler: No Free Lunch ve Inductive Bias}
\begin{itemize}
    \item \textbf{No Free Lunch (Bedava Öğle Yemeği Yok):} Matematiksel olarak, tüm olası problemler üzerinde diğerlerinden her zaman daha iyi performans gösteren tek bir "en iyi" algoritma yoktur. Model başarısı tamamen veri setinin yapısına ve probleme uygunluğuna bağlıdır.
    \item \textbf{Endüktif Önyargı (Inductive Bias):} Bir modelin eğitim verisinin ötesine geçip genelleme yapabilmesi için sahip olduğu varsayımlar bütünüdür. Örneğin, doğrusal regresyonun verideki ilişkinin doğrusal olduğunu varsayması bir önyargıdır.
\end{itemize}

\subsubsection{Makine Öğrenmesi İş Akışı (Pipeline)}
Sıfırdan başlayan bir projede takip edilen standart adımlar şöyledir:
\begin{enumerate}
    \item \textbf{Veri Toplama:} Ham verinin (CSV, SQL, API vb.) sisteme alınması.
    \item \textbf{Veri Ön İşleme:} Eksik değerlerin doldurulması ve ölçekleme (Bölüm 3).
    \item \textbf{Model Seçimi ve Eğitim:} Probleme uygun algoritmanın seçilip verilerle "eğitilmesi".
    \item \textbf{Değerlendirme:} Modelin hiç görmediği "test verisi" ile başarısının ölçülmesi.
    \item \textbf{Optimizasyon:} Hiperparametreleri değiştirerek en iyi sonucun aranması.
\end{enumerate}

\newpage
\subsection{Hata ve Performans Analizi}
Bir modelin başarısı, toplam hatanın bileşenlerine ayrıştırılmasıyla analiz edilir. Toplam hata (MSE), \textbf{İndirgenebilir Hata} ve \textbf{İndirgenemez Hata} ($\sigma^2$) olarak ikiye ayrılır.

\subsubsection{Bias-Variance Ödünleşimi (Trade-off)}
\begin{itemize}
    \item \textbf{Bias (Yanlılık):} Modelin verideki temel ilişkiyi kavrayamayacak kadar basit olması durumudur (Underfitting).
    \item \textbf{Variance (Varyans):} Modelin eğitim verisindeki gürültüyü ezberleyerek aşırı esnek hale gelmesi durumudur (Overfitting).
\end{itemize}
İdeal bir çalışma, her iki hatanın da dengelendiği noktayı bulmayı hedefler.

\subsubsection{Hata Bileşenleri: $MSE = Bias^2(f) + Var(f) + \sigma^2$}
Hata formülündeki $\sigma^2$, veri setindeki doğal gürültüdür ve hiçbir model tarafından yok edilemez. Mühendislik çabası, $Bias^2 + Var$ toplamını minimize etmeye yöneliktir.

\subsubsection{Sınıflandırma Hata Oranı ($E$)}
Sınıflandırma modellerinde temel başarı ölçütü olan hata oranı, tahmin edilen etiketin ($\hat{v}_i$) gerçek etiketten ($v_i$) farklı olduğu durumların toplam örneğe oranıdır:
\begin{equation}
    E = \frac{1}{n} \sum_{i=1}^{n} L(\hat{v}_i, v_i) \quad \text{burada } L=1 \text{ eğer } \hat{v}_i \neq v_i, \text{ aksi halde } 0
\end{equation}
Bu formül, dökümanın ilerleyen bölümlerinde detaylandırılacak olan "Confusion Matrix" (Karmaşıklık Matrisi) üzerindeki köşegen dışı elemanların toplamını temsil eder.

\subsubsection{Aşırı Öğrenme (Overfitting) ve Eksik Öğrenme (Underfitting)}
\begin{itemize}
    \item \textbf{Overfitting (Aşırı Uyum):} Eğitim hatası çok düşük, ancak test hatası çok yüksektir. Model verideki gürültüyü bile ezberlemiştir.
    \\ \textbf{Çözüm Stratejileri:} 
    \begin{itemize}
        \item Daha fazla veri toplamak.
        \item \textbf{Düzenlileştirme (Regularization):} Modelin katsayılarına ceza puanı vererek (L1/L2 gibi yöntemlerle) aşırı büyük değerler almasını engellemek.
        \item \textbf{Dropout:} Eğitim sırasında nöronların bir kısmını rastgele devre dışı bırakarak modelin belirli yollara bağımlı kalmasını önlemek.
        \item Model karmaşıklığını azaltmak (katman veya özellik sayısını düşürmek).
    \end{itemize}
    
    \item \textbf{Underfitting (Yetersiz Uyum):} Hem eğitim hem test hatası yüksektir. Model veriye dair temel yapıyı bile kavrayamamıştır.
    \\ \textbf{Çözüm Stratejileri:} 
    \begin{itemize}
        \item Daha karmaşık/derin modeller kullanmak.
        \item \textbf{Özellik Mühendisliği (Feature Engineering):} Yeni ve daha açıklayıcı özellikler türetmek.
        \item Düzenlileştirme (ceza) miktarını azaltmak.
    \end{itemize}
\end{itemize}

\subsection{Veri Yönetimi Stratejileri}
\subsubsection{Eğitim, Test ve Doğrulama (Validation) Setleri}
Modelin genelleme yeteneğini ölçmek için veri genellikle \%80 Eğitim, \%20 Test olarak bölünür. Hiperparametre optimizasyonu için ise Eğitim setinden bir "Doğrulama (Validation)" seti ayrılır.

\subsubsection{Model Seçimi: AIC ve BIC Kriterleri}
Hangi modelin daha iyi olduğunu belirlemek için hata payı ile parametre sayısını dengeleyen kriterlerdir:
\begin{itemize}
    \item \textbf{AIC (Akaike) ve BIC (Bayesian):} Değer ne kadar \textbf{küçükse} model o kadar iyidir. Bu kriterler, modele eklenen her gereksiz parametre için modeli cezalandırır.
\end{itemize}
\begin{itemize}
    \item \textbf{Parametre:} Modelin eğitim sırasında veriden kendi öğrendiği değerlerdir (Örn: Regresyondaki katsayılar/ağırlıklar).
    \item \textbf{Hiperparametre:} Eğitim başlamadan önce \textbf{bizim} dışarıdan belirlediğimiz ayar düğmeleridir (Örn: k-NN'deki $k$ sayısı, Sinir Ağlarındaki katman sayısı).
\end{itemize}

\subsubsection{Veri Sızıntısı (Data Leakage) Kavramı}
Test setindeki bilginin model eğitim aşamasında yanlışlıkla kullanılmasıdır. \textbf{Hayati Uyarı:} Bu durum, modelin gerçekte olduğundan çok daha başarılı görünmesine (sahte başarı) yol açar ve model gerçek veriye çıktığında çöker.




\newpage
\section{Veri Hazırlığı ve Özellik Mühendisliği}

\subsection{Veri Ön İşleme (Preprocessing)}
Ham verinin model tarafından işlenebilir "Sadık Embedding (Vektör)" yapılarına dönüştürülme sürecidir. Bir modelin başarısı, veriyi nasıl "gördüğüne" doğrudan bağlıdır.

\subsubsection{Ölçekleme: Standartlaştırma (Z-Score) ve Normalizasyon}
\begin{itemize}
    \item \textbf{Z-Skoru Standartlaştırma:} Verinin ortalamasını 0, standart sapmasını 1 yapar. Özelliklerin farklı ölçeklerde (Örn: Yaş ve Gelir) olmasından kaynaklanan "baskınlık" sorununu çözer.
    \item \textbf{Min-Max Normalizasyon:} Veriyi $[0, 1]$ aralığına sıkıştırır. Görüntü işleme gibi belirli sınırlar gerektiren alanlarda tercih edilir.
\end{itemize}

\subsubsection{Eksik Değer Analizi ve İmputasyon}
Veri setindeki boşluklar (NaN) model eğitimini durdurabilir.
\begin{itemize}
    \item \textbf{Temel Müdahaleler:} Satır silme (veri kaybı azsa) veya basit atama (ortalama/mod kullanımı).
    \item \textbf{Akıllı Tahminleme:} Eksik değerin, diğer değişkenler kullanılarak bir regresyon modeli veya PCA (Temel Bileşen Analizi) ile tahmin edilmesidir.
\end{itemize}

\subsubsection{Dağılım Yönetimi ve Dönüşümler}
\begin{itemize}
    \item \textbf{Log ve Gama Dönüşümü:} Web trafiği gibi çok geniş ölçekte değişen "çarpık" (skewed) verileri normale yaklaştırmak için kullanılır.
    \item \textbf{Winsorizing (Kırpma):} Aykırı değerlerin (outliers) etkisini azaltmak için veriyi belirli yüzdelik dilimler (Örn: \%5 ve \%95) arasında sabitleme işlemidir.
    \item \textbf{One-Hot Encoding:} Kategorik verileri (Elma, Armut, Muz) aralarında büyüklük ilişkisi kurmadan $[1, 0, 0]$ gibi birim vektörlere dönüştürür.
\end{itemize}

\subsection{Özellik Seçimi (Feature Selection)}
Gereksiz özelliklerin elenmesi; aşırı uyumu azaltır, hesap verimliliğini ve yorumlanabilirliği artırır.

\begin{enumerate}
    \item \textbf{Filtre Yöntemleri (Filter):} Modelden bağımsız, istatistiksel skorlara dayanır. Hızlıdır. (Örn: Pearson Korelasyonu, Chi-Kare, Mutual Information).
    \begin{itemize}
        \item \textbf{ANOVA F-Testi:} Sınıf ortalamaları arasındaki farkın, sınıf içi varyansa oranını ($F$ değeri) ölçer. Yüksek $F$ değeri, özelliğin ayırt edici gücünün yüksek olduğunu gösterir.
    \end{itemize}
    \item \textbf{Sarma Yöntemleri (Wrapper):} Özellik alt kümelerini bir model üzerinde test eder. Yavaştır ama doğrudur. (Örn: Forward Selection, RFE - Recursive Feature Elimination).
    \item \textbf{Gömülü Yöntemler (Embedded):} Özellik seçimi eğitimle birlikte gerçekleşir. (Örn: LASSO - L1 Düzenlileştirme, Karar Ağacı Önemi).
\end{enumerate}

\subsection{Boyutların Laneti (Curse of Dimensionality)}
Özellik sayısı arttıkça, veriyi temsil etmek için gereken örneklem miktarı üstel olarak artar. Bu durum geometrik bir seyrekliğe yol açar; yüksek boyutlarda veri noktaları birbirine çok uzak kalır ve benzerlik ölçümleri (mesafeler) anlamını yitirmeye başlar. Bu nedenle, modelin genelleme yeteneğini korumak için boyut düşürme veya özellik seçimi kritik önem taşır.

\subsection{Mesafe ve Benzerlik Ölçütleri}
Bir ölçümün matematiksel olarak \textbf{metrik} sayılabilmesi için üç temel kuralı sağlaması gerekir: Negatif olmama ($d(x,y) \geq 0$), simetri ($d(x,y)=d(y,x)$) ve \textbf{Üçgen Eşitsizliği} ($d(x,y) \leq d(x,z) + d(z,y)$).

\begin{itemize}
    \item \textbf{L1 (Manhattan) Mesafesi:} Izgara tabanlı şehir yapısındaki hareket gibi mutlak farkların toplamıdır. Yüksek boyutlu uzaylarda ve aykırı değerlere karşı Öklid'e göre daha dirençlidir.
    \item \textbf{L2 (Öklid) Mesafesi:} İki nokta arasındaki en kısa doğrusal mesafedir. Standart yaklaşımdır.
    \item \textbf{Kosinüs Benzerliği:} Vektörlerin büyüklüğüne değil, \textbf{yönüne ve aralarındaki açıya} odaklanır. Metin madenciliğinde (NLP) yaygındır.
    \item \textbf{Jaccard İndeksi:} Kümelerin ne kadar örtüştüğünü ölçer. İkili (binary) veriler için idealdir.
\end{itemize}



\newpage
\section{Regresyon Modelleri}
Regresyon; bağımsız değişkenler ($X$) ile sürekli bir bağımlı değişken ($y$) arasındaki ilişkiyi öğrenen denetimli bir öğrenme yöntemidir. Temel amaç, verilen bir girdi için sayısal bir tahmin üretmektir.

\subsection{Doğrusal ve Çoklu Regresyon}
\subsubsection{En Küçük Kareler (OLS) ve Normal Denklem}
OLS yöntemi, tahmin edilen değerler ile gerçek değerler arasındaki artıkların (residuals) kareleri toplamını (SSE) minimize etmeyi hedefler. Analitik çözüm olan Normal Denklem şu şekildedir:
\begin{equation}
    \beta = (X^T X)^{-1} X^T y
\end{equation}
\textbf{Güven Notu:} Bu formül işin matematiksel temelidir. Uygulamada bu matris işlemlerini Scikit-Learn gibi kütüphaneler otomatik olarak, milisaniyeler içinde yapar.
OLS, veri temiz ve az boyutlu olduğunda hızlıdır ancak aykırı değerlere (outliers) karşı oldukça duyarlıdır.

\subsubsection{Polinomsal Regresyon}
Doğrusal olmayan (eğrisel) ilişkileri modellemek için özellikleri yüksek dereceli terimlere ($x^2, x^3$) dönüştürerek doğrusal modellerle eğitme işlemidir. Aşırı karmaşık dereceler varyansı artırır.

\subsubsection{RANSAC: Sağlam (Robust) Regresyon}
Verideki aykırı değerlerin modeli bozmasını engellemek için kullanılır. Algoritma, veriden rastgele küçük bir alt küme (inliers) seçerek model kurar ve en çok veriyi açıklayan temiz modeli belirler.

\subsection{Düzenlileştirme (Regularization) ve Seyreklik}
Modelin katsayılarının aşırı büyümesini engelleyerek aşırı uyumu (overfitting) önleyen yöntemlerdir.

\begin{itemize}
    \item \textbf{Ridge Regresyon (L2):} Hata fonksiyonuna katsayıların karelerinin toplamını ekler ($\alpha \cdot \sum \beta_j^2$). Katsayıları sıfıra yaklaştırır ama tam sıfır yapmaz. Çoklu doğrusallık (multicollinearity) durumunda etkilidir.
    \item \textbf{Lasso Regresyon (L1):} Hata fonksiyonuna katsayıların mutlak değerlerinin toplamını ekler ($\alpha \cdot \sum |\beta_j|$). Bazı katsayıları tam olarak sıfır yaparak \textbf{Özellik Seçimi} sağlar. Seyrek (sparse) modeller üretir.
    \item \textbf{Elastic Net:} L1 ve L2 düzenlileştirmelerini birleştiren hibrit yaklaşımdır. Gruplu değişkenlerin olduğu senaryolarda Lasso'dan daha iyi performans gösterir.
\end{itemize}

\subsection{Hata Ölçütleri ve Yorumlama}
\begin{itemize}
    \item \textbf{MAE (Mean Absolute Error):} Hataların mutlak ortalamasıdır. Aykırı değerlere karşı stabildir.
    \item \textbf{MSE (Mean Squared Error):} Hataların karesinin ortalamasıdır. Büyük hataları çok daha sert cezalandırır. Türevi kolay alındığı için optimizasyonlarda sıkça tercih edilir.
    \item \textbf{RMSE (Root MSE):} MSE'nin kareköküdür. Hata birimini orijinal birime (Örn: TL, Derece) geri döndürerek yorumlamayı kolaylaştırır.
    \item \textbf{MAPE (Mean Absolute Percentage Error):} Hata payını yüzde olarak verir. Örn: "\%10 sapma ile çalışıyoruz". \textbf{Kritik Zayıflık:} Gerçek değerler sıfır veya sıfıra çok yakın olduğunda tanımsız hale gelir.
\end{itemize}

\subsubsection{R-Kare ve Düzeltilmiş R-Kare}
\textbf{R-Kare ($R^2$):} Bağımsız değişkenlerin, hedefteki varyansın yüzde kaçını açıkladığını gösterir.
\\ \textbf{Düzeltilmiş R-Kare ($Adjusted R^2$):} Modele eklenen gereksiz değişkenler için ceza uygular. Standart $R^2$ asla düşmezken, $Adjusted R^2$ model kalitesi artmadıkça yükselmez; böylece sahte başarının önüne geçilir.



\newpage
\section{Sınıflandırma Modelleri}
Sınıflandırma, veri noktalarını önceden tanımlanmış kategorilere (etiketlere) atama sürecidir. Çıktı değişkeni süreksiz ve kategoriktir.

\subsection{Lojistik Tahmin}
Lojistik regresyon, ismine rağmen bir sınıflandırma algoritmasıdır. Çıktıyı $[0, 1]$ arasına sıkıştırmak için \textbf{Sigmoid (Lojistik)} fonksiyonunu kullanır. İkili (binary) sınıflandırma problemlerinde (Örn: Hasta/Sağlıklı) temel yöntemdir.

\subsection{Örnek Tabanlı Öğrenme: k-NN}
k-En Yakın Komşu (k-NN) algoritması, \textbf{Tembel Öğrenme (Lazy Learning)} olarak bilinir. Eğitim aşamasında model kurmaz, tüm hesaplamayı tahmin anında yapar. Bir noktanın sınıfını, ona en yakın $k$ adet komşusunun oylarıyla belirler.
\\ \textbf{Kritik Detay:} $k$ sayısının çok küçük seçilmesi aşırı uyuma (overfitting), çok büyük seçilmesi ise yetersiz uyuma (underfitting) yol açar.

\subsection{Destek Vektör Makineleri (SVM)}
SVM, sınıflar arasındaki boşluğu (\textbf{Margin}) maksimize eden en uygun karar hiperdüzlemini bulmayı hedefler.
\begin{itemize}
    \item \textbf{Kernel Trick (Çekirdek Yöntemi):} Doğrusal olarak ayrılamayan verileri yüksek boyutlu bir uzaya taşıyarak ayrılabilir hale getirir.
    \item \textbf{Destek Vektörleri:} Karar sınırını belirleyen, sınırlara en yakın olan veri noktalarıdır.
\end{itemize}

\subsection{Olasılıksal Modeller: Naive Bayes}
Bayes Teoremi'ne dayanır. Özelliklerin sınıf verildiğinde birbirinden tamamen \textbf{bağımsız} olduğunu varsaydığı için "Naive" (Saf) olarak adlandırılır. Hızlıdır ve özellikle metin sınıflandırma (Spam filtresi) işlerinde çok etkilidir.
\\ \textbf{Laplace Smoothing:} Veride olmayan bir özelliğin olasılığının 0 olmasını ve tüm sonucu sıfırlamasını engellemek için tüm sayımlara 1 eklenmesidir.

\subsection{Karar Yapıları ve Topluluklar (Ensembles)}
\begin{itemize}
    \item \textbf{Karar Ağaçları:} Veriyi \textbf{Entropi} veya \textbf{Gini Safsızlığı} ölçütlerine göre dallara ayırarak karar verir. Okunabilirliği yüksektir.
    \item \textbf{Random Forest (Bagging):} Çok sayıda karar ağacını rastgele örneklemle (bootstrap) eğitip sonuçları oylayarak varyansı düşürür.
    \item \textbf{Gradient Boosting (Boosting):} Ağaçları ardışıl (sequential) olarak eğitir; her yeni ağaç bir önceki ağacın hatalarını düzeltmeye odaklanır.
\end{itemize}


\newpage
\section{Kümeleme ve Yapı Bulma}
Kümeleme, veri setindeki benzer örnekleri gruplandırmayı amaçlayan denetimsiz bir öğrenme yöntemidir. Önceden tanımlanmış bir etiket (label) yoktur; yapı verinin kendi özelliklerinden keşfedilir.

\subsection{Merkez Tabanlı Kümeleme: K-Means}
Veriyi $K$ adet kümeye ayıran iteratif bir algoritmadır. Her nokta en yakın küme merkezine (centroid) atanır ve merkezler her adımda güncellenir.
\begin{itemize}
    \item \textbf{Inertia:} Veri noktalarının kendi küme merkezlerine olan uzaklıklarının kareleri toplamıdır. Düşük inertia, daha sıkı kümeler anlamına gelir.
    \item \textbf{Dirsek (Elbow) Yöntemi:} Optimal küme sayısını ($K$) belirlemek için kullanılır. Inertia'daki düşüşün "dirsek" yaptığı nokta seçilir.
    \item \textbf{Silhouette Skoru:} Bir noktanın kendi kümesine ne kadar benzer, komşu kümelere ne kadar uzak olduğunu ölçer. $+1$ mükemmel kümelemeyi, $-1$ hatalı atamayı gösterir.
\end{itemize}

\subsection{Hiyerarşik ve Yoğunluk Tabanlı Yapılar}
\begin{itemize}
    \item \textbf{Hiyerarşik Kümeleme:} Veriyi ağaç benzeri bir yapıda (\textbf{Dendrogram}) birleştirir. \textbf{Single Linkage} (en yakın komşu) veya \textbf{Complete Linkage} (en uzak komşu) gibi mesafe kuralları kullanılır.
    \item \textbf{DBSCAN:} Yoğunluk tabanlı bir yöntemdir. Karmaşık şekilli kümeleri bulabilir ve \textbf{Gürültüyü} (aykırı değerleri) diğer kümelerden ayırabilir.
\end{itemize}

\subsection{İleri Seviye Temsil}
\textbf{Gauss Karışım Modelleri (GMM):} Verinin farklı olasılık dağılımlarının birleşimi olduğunu varsayar ve "yumuşak" (soft) kümeleme yapar. \textbf{EM Algoritması} ile optimize edilir.


\newpage
\section{Boyut Düşürme ve Veri Sıkıştırma}
Yüksek boyutlu verilerdeki gereksiz varyansı temizleyerek veriyi daha düşük boyutlu bir alt uzaya izdüşürme işlemidir.

\subsection{Doğrusal İzdüşüm Yöntemleri}
\begin{itemize}
    \item \textbf{PCA (Temel Bileşen Analizi):} Verideki maksimum varyansı koruyarak birbirine dik (\textbf{Ortogonal}) yeni bileşenler üretir. Bu süreçte kovaryans matrisinin \textbf{Özvektörleri (Eigenvectors)} ve önem derecelerini belirleyen \textbf{Özdeğerleri (Eigenvalues)} kullanılır.
    \item \textbf{LDA (Doğrusal Diskriminant Analizi):} PCA'nın aksine denetimli bir yöntemdir. Sınıf ayrımını (separability) maksimize etmeyi hedefler. Sınıf sayısı $C$ ise, en fazla $C-1$ adet yeni bileşen üretebilir.
\end{itemize}

\subsection{Karmaşık ve Doğrusal Olmayan Dönüşümler}
\begin{itemize}
    \item \textbf{Otoenkodörler (Autoencoders):} Veriyi \textbf{Bottleneck} (Darboğaz) katmanında sıkıştırıp tekrar açan sinir ağlarıdır. Gürültü temizleme (\textbf{Denoising}) için de kullanılır.
    \item \textbf{t-SNE ve UMAP:} Özellikle verinin 2 veya 3 boyutta görselleştirilmesi için kullanılan, komşuluk ilişkilerini koruyan algoritmalardır.
\end{itemize}


\newpage
\section{Yapay Sinir Ağları ve Derin Öğrenme}
İnsan beynindeki nöron yapısından esinlenen, karmaşık doğrusal olmayan ilişkileri öğrenebilen katmanlı yapılardır.

\subsection{Nöron Yapısı ve Aktivasyon Fonksiyonları}
Bir yapay nöron; girdileri (\textbf{Weights}), bir yanlılık değerini (\textbf{Bias}) alır ve bunları bir aktivasyon fonksiyonundan geçirerek çıktı üretir.
\begin{itemize}
    \item \textbf{ReLU:} Negatif değerleri sıfırlar, doğrusal olmayan yapıyı hızlı şekilde öğrenir.
    \item \textbf{Sigmoid/Softmax:} Çıktıyı olasılık değerine dönüştürür.
\end{itemize}

\subsection{Eğitim Süreci: Geri Yayılım (Backpropagation)}
Çıkış katmanında oluşan hatanın, Zincir Kuralı (Chain Rule) kullanılarak geriye doğru ağa dağıtılması ve ağırlıkların bu hataya göre güncellenmesidir. Bu süreç, modelin tahmindeki payı olan ağırlıkları kademeli olarak optimize etmesini sağlar.

\subsection{Optimizasyon ve Düzenlileştirme}
Derin ağlarda \textbf{Gradyan Kaybolması} (Vanishing Gradient) sorununu aşmak için \textbf{Batch Normalization} ve \textbf{Dropout} gibi teknikler kullanılır. \textbf{Evrensel Yaklaşım Teoremi}'ne göre, yeterince geniş bir ağ her türlü sürekli fonksiyonu temsil edebilir.


\newpage
\section{Model Değerlendirme ve Optimizasyon}

\subsection{Doğrulama ve Çapraz Doğrulama (Cross-Validation)}
\begin{itemize}
    \item \textbf{k-Fold CV:} Veriyi $k$ parçaya böler. Genellikle $k=5$ veya $k=10$ seçilir.
    \item \textbf{Bootstrap (Önyükleme):} Veriden \textbf{yerine koyarak} (with replacement) örnek çekilmesidir.
    \item \textbf{LOOCV (Leave-One-Out):} Küçük verilerde hassastır ancak yüksek hesaplama maliyeti dezavantajıdır.
    \item \textbf{Nested (İç İçe) CV:} Aynı anda hem hiperparametre ayarlamak hem de tarafsız performans ölçmek için kullanılır.
\end{itemize}

\subsection{Hiperparametre Optimizasyonu Stratejileri}
\begin{itemize}
    \item \textbf{Grid Search (Izgara Araması):} Belirlenen tüm kombinasyonları dener. Hesaplama maliyeti $O(L^F)$ formülüyle ifade edilir ($L$: Seviye, $F$: Faktör).
    \item \textbf{Random Search (Rastgele Arama):} Yüksek boyutlu uzaylarda, önemli olan parametreleri daha verimli taradığı için Grid Search'ten daha başarılı olabilir.
    \item \textbf{Simulated Annealing:} "Soğuma" dinamiğiyle yerel minimumlardan kaçmayı hedefler.
    \item \textbf{Genetik Algoritmalar:} \textbf{Mutasyon ve Çaprazlama} mekanizmalarıyla en iyi parametreleri arar.
\end{itemize}

\subsection{Başarı Metrikleri}
\begin{itemize}
    \item \textbf{Karmaşıklık Matrisi:} TP, TN, FP (Tip I Hata), FN (Tip II Hata) değerlerini gösterir.
    \item \textbf{F1 Skoru:} Precision ve Recall'un \textbf{Harmonik Ortalamasıdır}. Harmonik ortalama, uç (çok düşük) değerleri şiddetle cezalandırarak dengeli bir başarı ister.
    \item \textbf{ROC-AUC:} Modelin sınıfları ayırma kabiliyetini [0, 1] aralığında ölçer. $AUC=0.5$ rastgele tahmini, $AUC=1$ mükemmel ayrımı ifade eder.
    \item \textbf{Neyman-Pearson Kriteri:} Hata maliyetlerinin farklı olduğu durumlarda kullanılır. Toplam hata $E = \lambda \frac{C_{fn}}{n} + \frac{C_{fp}}{n}$ formülüyle hesaplanır. Burada $\lambda > 1$ seçilmesi, yanlış negatiflerin (vaka kaçırmanın) daha maliyetli kabul edildiği anlamına gelir.
\end{itemize}


\newpage
\section{Özet Karşılaştırmalar ve Kritik İpuçları}
Bu bölüm, öğrenme kartlarında ve sınav bankasında sıkça karşınıza çıkabilecek "ayırt edici" detayları özetlemektedir:

\subsection{Regresyon ve Düzenlileştirme}
\begin{itemize}
    \item \textbf{Ridge (L2) vs. Lasso (L1):} Lasso katsayıları tam olarak sıfıra çekebilir ve otomatik özellik seçimi yapar. Ridge ise katsayıları küçültür ama sıfırlamaz.
    \item \textbf{R-Kare (R2):} Modelin veriyi açıklama oranını gösterir. Ayarlanmış R2 ise modele eklenen gereksiz özelliklerin başarısını cezalandırır.
\end{itemize}

\subsection{Mesafe ve Sınıflandırma}
\begin{itemize}
    \item \textbf{Mesafe Seçimi:} Sayısal verilerde \textbf{Euclidean (L2)}, aykırı değerlere hassas durumlarda \textbf{Manhattan (L1)}, metin veya açısal verilerde \textbf{Cosine}, kümeler arası benzerlikte \textbf{Jaccard} kullanılır.
    \item \textbf{Parametre Sayısı (LDA):} LDA, $C$ adet sınıfın olduğu bir problemde en fazla $C-1$ adet yeni boyut (bileşen) üretebilir.
\end{itemize}

\subsection{Karar Ağaçları ve Topluluklar}
\begin{itemize}
    \item \textbf{Safsızlık Ölçütleri:} \textbf{Entropi} (bilgi kazancı odaklı) ve \textbf{Gini Safsızlığı} (daha hızlı hesaplanır, Scikit-Learn varsayılanıdır) ağacı bölmek için kullanılır.
    \item \textbf{Bagging vs. Boosting:} Bagging (Random Forest) varyansı düşürür; Boosting (Gradient Boosting) hem bias hem varyansı düşürerek daha karmaşık hataları çözer.
\end{itemize}

\newpage
\section*{Kaynakça ve Teşekkür}
\addcontentsline{toc}{section}{Kaynakça ve Teşekkür}
Bu çalışma, İstanbul Medeniyet Üniversitesi, Mühendislik ve Doğa Bilimleri Fakültesi bünyesinde verilen \textbf{Makine Öğrenmesine Giriş} dersi kapsamında hazırlanan akademik materyallerden sentezlenmiştir.

Dökümanın hazırlanmasında kullanılan temel teorik bilgiler, sunumlar ve ders notları; değerli hocamız \textbf{Doç. Dr. Betül Hiçdurmaz} tarafından paylaşılan kaynaklara dayanmaktadır. Akademik katkıları ve paylaştığı zengin içerikler için hocamıza teşekkürlerimizi sunarız.

\end{document}
